\newpage
\section{问题分析}

\subsection{问题一分析}
本题要求预测碳排放趋势并识别碳达峰时间，属于\textbf{时序预测+转折点检测}
问题。Prophet模型擅长处理具有强季节性和节假日效应的时序数据。我们需要确定最优的changepoint_prior_scale参数，使模型既能捕捉长期趋势，又不拟合噪声。Prophet模型形式为：\y(t) = g(t) + s(t) + h(t) + \epsilon_t
其中g(t)为趋势项（分段线性或Logistic增长），s(t)为周期项（傅里叶级数），h(t)为节假日效应。

\subsection{问题二分析}
本题是\textbf{因素分解}
问题。LMDI方法将碳排放变化分解为多个驱动因子：\C = P \cdot (GDP/P) \cdot (E/GDP) \cdot (C/E)
分别对应人口、人均GDP、能源强度和碳排放强度。LMDI分解具有完全分解性（无残差项），适合用于政策归因分析。对数平均权重函数为：\L(a,b) = \frac{a-b}{\ln a - \ln b}, \quad a \neq b

\subsection{问题三分析}
本题是\textbf{多目标优化}
问题。三个目标之间存在内在冲突：降低碳排放需要增加清洁能源投资，推高经济成本。采用加权求和法将多目标转化为单目标：\\min \quad Z = w_1 \cdot \frac{C_{cost}}{C_{cost}^{max}} + w_2 \cdot \frac{E_{sec}}{E_{sec}^{max}} + w_3 \cdot \frac{C_{carbon}}{C_{carbon}^{max}}
权重通过熵权法客观确定，最终得到Pareto最优解集。